\subsection{The remaining regular-representability gate}
The August 5, 2026 revision of
\cite[Introduction, p.~1]{DGJ_ACMVTwist2026} still distinguishes
universal algebraic representatives in codimensions $1,2,\dim X$
from merely regular algebraic intermediate-Jacobian maps in
the other codimensions. This is a precise current source check,
not an inference of open status from an unsuccessful search.
The final published version of Kahn's concrete fourfold
expectation is now also opened: Conjecture~3.1, p.~847
\cite{DGJ_Kahn2021}. It corresponds to Conjecture~1, p.~5 in
the six-page author version cited above.

For complex varieties one should additionally distinguish
universality among all regular homomorphisms from universality
among isogeny lifts. Write
\[
 \psi:A^i(X)\longrightarrow J_a^{2i-1}(X)(\mathbb C),\qquad
 \omega:A^i(X)\longrightarrow J_W^{2i-1}(X)(\mathbb C)
\]
for the usual and Walker Abel--Jacobi maps, and
$\alpha:J_W^{2i-1}(X)\to J_a^{2i-1}(X)$ for their isogeny.
By \cite[Theorem 2.3, p.~12]{DGJ_ACMVWalker}, $\omega$ is
initial among regular lifts of $\psi$ along isogenies. This
theorem does not say that $\omega$ is universal among all
regular homomorphisms.

Here is an exact elementary test for that stronger assertion.
For every regular homomorphism
$\phi:A^i(X)\to B(\mathbb C)$, let
\[
 C_\phi=\operatorname{im}(\psi,\phi)
       \subset J_a^{2i-1}(X)\times B.
\]
The image is the group of complex points of an abelian
subvariety, by the image theorem for regular homomorphisms
\cite[Lemma 4.7, pp.~18--19]{DGJ_ACMVFunctor}.
Then $\omega$ is universal among all regular homomorphisms
if and only if
\begin{equation}\label{eq:JWalker-gate}
 \dim C_\phi=\dim J_a^{2i-1}(X)
 \quad\hbox{for every regular }\phi.
\end{equation}
Indeed, if \eqref{eq:JWalker-gate} holds, the surjective
projection $C_\phi\to J_a^{2i-1}(X)$ is an isogeny. Initiality
among isogeny lifts gives a homomorphism
$J_W^{2i-1}(X)\to C_\phi$ inducing $(\psi,\phi)$; its second
projection factors $\phi$ through $\omega$. This factorization
is unique because $\omega$ is surjective. Conversely, if
$\omega$ is universal, $(\psi,\phi)$ factors through $J_W$,
so its image has dimension at most $\dim J_W=\dim J_a$;
the surjective projection to $J_a$ gives the reverse inequality.

The unresolved issue in this route is therefore a possible
positive-dimensional kernel of $C_\phi\to J_a$ for a regular
test $\phi$. The divisible-kernel statement in Theorem~2.3
does not rule it out: divisible groups can have nonzero maps
to groups of points of abelian varieties. Likewise, infinite
torsion in $A^3$ does not by itself construct regular targets
of unbounded dimension. The test above is not asserted for
Kahn's fourfold, and neither its proof nor its failure has
been obtained here. This leaves the general statement in
Theorems~1--3 explicitly unresolved.
